% Section 3.4, from lectures/L03/appendix/b_exercises.md. The solutions to the three hand-training
% examples, which the appendix gives inline, are in Appendix A; the solution to the code is in the
% repository.

\section{Exercises}
\label{c3:sec:exercises}\label{c3:app:b}

\begin{aside}
\textbf{How to work these.} Set~1 is three calculations: work them on paper with the equations
from \secref{c3:sec:nn} and a calculator, one training set at a time. Set~2 starts a new \code{ml}
codebase in this lecture's exercises directory, \file{lectures/L03/appendix/exercises} in the
companion repository, and is built with a Makefile you write yourself.

\textbf{How to check your work.} The worked solutions to Set~1 are in
Appendix~\ref{sol:ch:solutions}. Set~2 comes with a test suite, already in
\file{lectures/L03/appendix/exercises/test}, which Exercise~\ref{c3:ex:7} runs. It tests the
directory above it by default; \code{make ML_DIR=<path>} points it at a tree of your own elsewhere.

The solution to Set~2 is published in the companion repository after the lecture;
\file{lectures/L03/README.md} links to it. Try each exercise yourself before looking at it.
\end{aside}

% ------------------------------------------------------------------------------------------
\exerciseset{Hand-Training a Neural Network}
\label{c3:set:1}

\calcexercise{Odd parity, first training set}
\label{c3:ex:1}
The neural network in Figure~\ref{fig:network1} should be trained to predict a high output $Y$
when an odd number of the inputs $X_1$ and $X_2$ are high, and output $Y = 0$ otherwise:

\bookfigure{network1}{The network in Exercises~\ref{c3:ex:1} and~\ref{c3:ex:3}: two inputs, two
hidden nodes and one output node.}{fig:network1}

\noindent The first training set: inputs $X_1X_2 = 01$, $Y = 1$.

\task{Starting parameters}
\begin{align*}
  b_1 &= 0.2, & b_2 &= 0.4, & b_3 &= 0.6 \\
  w_1 &= 0.5, & w_2 &= 0.6, & w_3 &= 0.3 \\
  w_4 &= 0.8, & w_5 &= 0.1, & w_6 &= 0.9
\end{align*}
Learning rate: $\LR = 0.1$. Activation function: ReLU for all nodes.

Train the network on the first training set: perform feedforward, backpropagation, and
optimization. Then feed the same input forward again with the new parameters, and compare the
deviation against the one you started with.
\seesolution{c3:ex:1}

\calcexercise{Inverting two inputs}
\label{c3:ex:2}
Consider the neural network in Figure~\ref{fig:network2}:

\bookfigure{network2}{The network in Exercise~\ref{c3:ex:2}: two inputs, three hidden nodes and
two output nodes.}{fig:network2}

\noindent The network should be trained to predict outputs $Y_1Y_2$ that form the inverse of inputs
$X_1X_2$:

\begin{narrowtable}{@{}cc@{}}
  $X_1X_2$ & $Y_1Y_2$ \\
  \midrule
  00 & 11 \\
  01 & 10 \\
  10 & 01 \\
  11 & 00 \\
\end{narrowtable}

\noindent First training set: $X_1X_2 = 10$, $Y_1Y_2 = 01$.

\task{Starting parameters: hidden layer}
\begin{align*}
  b_1 &= 0.1, & b_2 &= 0.1, & b_3 &= 0.7 \\
  w_1 &= 0.2, & w_2 &= 0.9, & w_3 &= 0.3 \\
  w_4 &= 0.8, & w_5 &= 0.5, & w_6 &= 0.4
\end{align*}

\task{Starting parameters: output layer}
\begin{align*}
  b_4 &= 0.1, & b_5 &= 0.6 \\
  w_7 &= 0.1, & w_8 &= 0.1, & w_9 &= 1.0 \\
  w_{10} &= 0.3, & w_{11} &= 0.0, & w_{12} &= 0.6
\end{align*}
Learning rate: $\LR = 0.1$. Activation function: ReLU for all nodes.

Train the network on the first training set: perform feedforward, backpropagation, and
optimization. Then feed the same input forward again with the new parameters, and compare the
deviations against the ones you started with.
\seesolution{c3:ex:2}

\calcexercise{Odd parity, after one epoch}
\label{c3:ex:3}
Consider the neural network in Figure~\ref{fig:network1} again. The network should be trained to
predict a high output $Y$ when an odd number of the inputs $X_1$ and $X_2$ are high, and output
$Y = 0$ otherwise.

Training has already been carried out for one epoch, and the parameters have the following starting
values:

\task{Starting parameters}
\begin{align*}
  b_1 &= 0.1, & b_2 &= -0.4, & b_3 &= -0.2 \\
  w_1 &= 0.5, & w_2 &= 0.4, & w_3 &= -0.2 \\
  w_4 &= 0.1, & w_5 &= 0.7, & w_6 &= 0.8
\end{align*}
Training is now to be carried out with the training set $X_1X_2 = 11$, $Y = 0$.

Learning rate: $\LR = 0.1$. Activation function: ReLU for all nodes.

Perform feedforward, backpropagation, and optimization. Then compute the output again. Did the
error decrease? If not, what could you have changed to obtain a smaller error?
\seesolution{c3:ex:3}

% ------------------------------------------------------------------------------------------
\exerciseset{Dense Layer Interface and Stub}
\label{c3:set:2}
You'll start a new \code{ml} codebase (or extend your existing one, if you'd like to keep the linear
regression code alongside it) with an interface for a dense layer and a placeholder implementation
of it.

\codeexercise{Directory structure}
\label{c3:ex:4}
Set up the following structure in this lecture's exercises directory,
\file{lectures/L03/appendix/exercises}:

\begin{console}
exercises/
├── include/
│   └── ml/
│       ├── dense_layer/
│       │   ├── interface.hpp
│       │   └── stub.hpp
│       └── types.hpp
├── source/
│   └── main.cpp
├── test/                <- already here, see Exercise 3.7
└── Makefile
\end{console}

\noindent You'll extend this structure further in Chapter~\ref{c4:ch:network}.

\codeexercise{Dense layer interface and stub class}
\label{c3:ex:5}
The hidden layer and output layer of a future neural network will be represented by the interface
\code{ml::dense_layer::Interface}. A concrete implementation isn't created until
Chapter~\ref{c5:ch:dense}. Until then, you'll implement a simple stub class
\code{ml::dense_layer::Stub} as a placeholder.

\task{The interface (\file{ml/dense_layer/interface.hpp})}
In the namespace \code{ml::dense_layer}, implement an interface named \code{Interface}. All methods
(except the destructor) should be declared pure virtual (\code{= 0}).
\begin{itemize}
  \item \textbf{\code{~Interface()}:} should be set to \code{default} and marked \code{virtual} and
    \code{noexcept}.
\end{itemize}
Getters, all \code{const}, \code{noexcept}, and \code{[[nodiscard]]}:

\begin{booktable}{@{}p{7.2em}L@{}}
  \thead{Method} & \thead{Returns} \\
  \midrule
  \code{nodeCount()} & Number of nodes in the layer (\code{std::size_t}). \\
  \code{weightCount()} & Number of weights per node (\code{std::size_t}). \\
  \code{output()} & Reference to the layer's output (read-only floating-point vector). \\
  \code{error()} & Reference to the layer's error (read-only floating-point vector). \\
  \code{weights()} & Reference to the layer's weights (read-only, two-dimensional floating-point
    vector). \\
\end{booktable}

\noindent Computation methods, all \code{noexcept} and returning \code{bool}. Each returns
\code{false} on invalid input (wrong dimensions or an invalid learning rate) and \code{true}
otherwise, so the caller can decide what to do about the failure. As in
Chapters~\ref{c1:ch:linreg} and~\ref{c2:ch:training}, \code{std::terminate()} is reserved for the
constructor, which has no way to return a failure code to the caller:
\begin{itemize}
  \item \textbf{\code{feedforward(input)}:} performs feedforward.
  \begin{itemize}
    \item \code{input}: read-only floating-point vector of input data.
  \end{itemize}
  \item \textbf{\code{backpropagate(reference)}} (output layer): computes error from reference
    values.
  \begin{itemize}
    \item \code{reference}: read-only floating-point vector of reference values.
  \end{itemize}
  \item \textbf{\code{backpropagate(nextLayer)}} (hidden layer): computes error from the next
    layer.
  \begin{itemize}
    \item \code{nextLayer}: reference to the next layer (\code{const Interface&}).
  \end{itemize}
  \item \textbf{\code{optimize(input, learningRate)}:} updates bias and weights.
  \begin{itemize}
    \item \code{input}: read-only floating-point vector.
    \item \code{learningRate}: floating-point number.
  \end{itemize}
\end{itemize}
One method more, \code{noexcept} and returning nothing:
\begin{itemize}
  \item \textbf{\code{initParams()}:} resets the layer's trainable parameters, i.e.\ its bias
    values and weights, to their initial values.
  \begin{itemize}
    \item Takes no arguments and returns nothing.
    \item Pure virtual like the rest, so every layer states what its own reset means, even when the
      answer is ``nothing''. \code{Stub} has no trainable parameters to draw again, so its override
      is empty; \code{Dense} overrides it in Chapter~\ref{c5:ch:dense} by drawing a new random bias
      and weight for every node, and its constructor calls it too, so the randomization is written
      once and used from both places.
  \end{itemize}
\end{itemize}

\textbf{Who calls it.} The network you write in Chapter~\ref{c4:ch:network} does, on both of its
layers, before its first training epoch. That's what makes a second call to \code{train()} train a
new network rather than continue the one the previous call left behind, which matters when a run
ends up badly trained: a network that has settled into a bad set of parameters mostly stays there,
while a fresh start draws new ones and can converge on the next attempt. The bias and weights are
private to the layer, so without a method on the interface the network couldn't reset them at all.

\task{The stub class (\file{ml/dense_layer/stub.hpp})}
In the namespace \code{ml::dense_layer}, implement a subclass named \code{Stub} that inherits
\code{Interface} via public inheritance. The class should be marked \code{final}. The stub doesn't
perform any real computation; it exists solely so other code can be compiled, test-run, and unit
tested against a real \code{dense_layer::Interface} before a concrete \code{Dense} implementation
exists (see Chapter~\ref{c5:ch:dense}). The network you build in Chapter~\ref{c4:ch:network} is
tested entirely against this stub, so it's worth getting right.
\begin{itemize}
  \item \textbf{\code{Stub()}:}
  \begin{itemize}
    \item The class's only implemented constructor.
    \item Should take the following arguments:
    \begin{itemize}
      \item \code{nodeCount}: number of nodes in the layer (unsigned integer).
      \item \code{weightCount}: number of weights per node (unsigned integer).
      \item \code{outputValue}: value every output element is set to (floating-point number).
        Default value: \code{0.5}.
    \end{itemize}
    \item Initializes every element of the output vector to \code{outputValue}, and the error and
      weight vectors to zeros.
    \item Should print an error message and call \code{std::terminate()} if either
      \code{nodeCount} or \code{weightCount} is zero:
    \begin{itemize}
      \item As in Chapters~\ref{c1:ch:linreg} and~\ref{c2:ch:training}, the constructor is the one
        place that terminates, since it has no way to return a failure code.
      \item It also guarantees at least one node and one weight, which is what makes it safe for
        \code{weightCount()} to read the width of the first row of the weight matrix.
    \end{itemize}
    \item \code{outputValue} is an argument rather than a hard-coded constant so that two stubs can
      be told apart. A network whose layers all report the same output can't show whether a
      prediction came from the output layer or from the hidden one; giving each layer its own value
      makes that visible.
    \item Should be marked \code{explicit} and \code{noexcept}.
  \end{itemize}
  \item \textbf{\code{~Stub()}:}
  \begin{itemize}
    \item Destructor overriding the interface's destructor.
    \item Should be marked \code{default}, \code{noexcept}, and \code{override}.
  \end{itemize}
  \item \textbf{Getters} (\code{nodeCount()}, \code{weightCount()}, \code{output()},
    \code{error()}, \code{weights()}):
  \begin{itemize}
    \item Override the corresponding methods in the interface.
    \item Should be marked \code{override}, retaining the interface's \code{const} and
      \code{noexcept}, but \textbf{not} \code{[[nodiscard]]}.
  \end{itemize}
  \item \textbf{\code{feedforward()}}, both overloads of \textbf{\code{backpropagate()}}, and
    \textbf{\code{optimize()}}:
  \begin{itemize}
    \item Perform range checks only:
    \begin{itemize}
      \item Return \code{false} when the dimensions don't match, or, for \code{optimize()}, when
        the learning rate lies outside \code{(0.0, 1.0)}.
      \item Return \code{true} otherwise.
    \end{itemize}
    \item Deliberately compute nothing:
    \begin{itemize}
      \item The output stays at \code{outputValue} regardless of what's fed in.
      \item The error stays at zero.
    \end{itemize}
    \item Should be marked \code{override} and \code{noexcept}.
  \end{itemize}
  \item \textbf{\code{setOutput()}:}
  \begin{itemize}
    \item Sets every element of the output vector to the given value.
    \item Should take a single argument:
    \begin{itemize}
      \item \code{outputValue}: value to set every output element to (floating-point number).
    \end{itemize}
    \item Returns nothing, and should be marked \code{noexcept}.
    \item \textbf{Not} part of \code{Interface}. It exists on the stub alone.
    \item It's what lets a test drive the output of a whole network in
      Chapter~\ref{c4:ch:network}:
    \begin{itemize}
      \item The network stores its layers by reference, so calling \code{setOutput()} on the layer
        a network was built with changes what that network predicts.
      \item That reveals whether the network reads its output layer live or kept a stale copy of
        it. Keeping a stale copy is exactly what the note in Chapter~\ref{c4:ch:network} about not
        needing a separate storage variable is there to prevent.
    \end{itemize}
  \end{itemize}
  \item \textbf{\code{feedforwardCount()}} and \textbf{\code{clearFeedforwardCount()}}:
  \begin{itemize}
    \item \code{feedforwardCount()} returns the number of times \code{feedforward()} has been
      called on this layer as a \code{std::size_t}, counting \textbf{every} call, not just the ones
      that passed the range check. Return \code{std::size_t} rather than \code{int}: the tests
      compare the tally against \code{std::size_t} values, and a signed return trips
      \code{-Wsign-compare}, which the suite builds with \code{-Werror}.
    \begin{itemize}
      \item Increment the counter (\code{myFeedforwardCount}) at the very top of
        \code{feedforward()}, before the input size is checked, so a rejected call raises it just
        as an accepted one does.
      \item The tally records how often the layer was \emph{asked} to feed forward, not how often
        it agreed to.
      \item Should be marked \code{[[nodiscard]]}, \code{const}, and \code{noexcept}.
    \end{itemize}
    \item \code{clearFeedforwardCount()} resets that tally to zero.
    \begin{itemize}
      \item Returns nothing, and should be marked \code{noexcept}.
    \end{itemize}
    \item Neither is part of \code{Interface}.
    \item \code{train()} in Chapter~\ref{c4:ch:network} performs one feedforward per training set
      per epoch, so this is what lets a test pin down the training loop:
    \begin{itemize}
      \item Nothing else can: a loop that runs a single pass instead of every epoch still lines up
        dimensionally and still returns \code{true}, so without a count it's indistinguishable from
        a correct one.
    \end{itemize}
  \end{itemize}
  \item \textbf{\code{initParams()}:}
  \begin{itemize}
    \item Overrides the corresponding method in the interface, and should be marked
      \code{override} and \code{noexcept}.
    \item Its body is empty: the stub has no bias, and its weights never move from the zeros it was
      constructed with, so there's nothing to draw again.
    \item An empty body is the implementation here, not a placeholder. It's also the only one the
      network in Chapter~\ref{c4:ch:network} needs: that network calls \code{initParams()} on both
      of its layers before every training run, and against a stub the call has to compile, run, and
      change nothing.
  \end{itemize}
\end{itemize}
For this class, the default constructor and the copy and move constructors (and corresponding
operators) should be deleted.

Add private member variables for the output, the error, the weights, and the feedforward count.

The number of nodes and the number of weights per node don't need member variables of their own.
Both are readable from the vectors you already have: \code{nodeCount()} is the size of the output
vector, and \code{weightCount()} is the width of the weight matrix.

The bias doesn't need one at all. It isn't part of \code{Interface}, and this stub never optimizes
anything, so there'd be nothing to put in it and no way to read it back out.

\note[Note!] Leaving it out is a stub-only simplification. The real \code{Dense} layer you write in
Chapter~\ref{c5:ch:dense} does need a bias vector: it's a trainable parameter, adjusted alongside the
weights on every call to \code{optimize()}.

\codeexercise{A quick compile check}
\label{c3:ex:6}
There's no network to run yet; that's Chapter~\ref{c4:ch:network}, once
\code{neural_network::Shallow} exists to make use of this interface. For now, just confirm your
\code{Interface}/\code{Stub} pair compiles: in \file{main.cpp}, create a \code{ml::dense_layer::Stub}
with a few nodes and weights, and print its \code{nodeCount()}, \code{weightCount()}, and
\code{output()} to the terminal. You should see \code{nodeCount()}/\code{weightCount()} match what
you constructed it with, and every value in \code{output()} equal to the default \code{outputValue}
of \code{0.5}.

Create a second stub with an explicit \code{outputValue}, and check that its \code{output()} reports
that value instead. Also feed both stubs a vector of the wrong size and check that
\code{feedforward()} returns \code{false} rather than terminating.

\codeexercise{Running the tests}
\label{c3:ex:7}
A test suite for the stub is available in \file{lectures/L03/appendix/exercises/test}. It's already
in place next to the code you just wrote. Build and run it from the exercises directory:

\begin{shell}
make -C test
\end{shell}

\noindent The test framework is checked out as a submodule at the root of the repository. If
\file{libs/test} is empty, fetch it once with \code{git submodule update --init}.

All 14 test cases should pass. \textbf{This chapter starts a new \code{ml} codebase}, so this suite
replaces nothing and stands on its own.

Testing a class that computes nothing may look pointless, but the range checks are the only real
logic the stub has, and the network you write in Chapter~\ref{c4:ch:network} is tested entirely
against them. A stub whose \code{feedforward()} accepted an input of any size would let every test
in Chapter~\ref{c4:ch:network} pass without proving anything.

The test suite's README, \file{lectures/L03/appendix/exercises/test/README.md}, has more
information, including which of these tests carry over unchanged to your real \code{Dense} layer in
Chapter~\ref{c5:ch:dense}.
